\begin{table}[ht]
\centering
\small
\begin{tabular}{lrllrr}
\toprule
\textbf{Population} & $n$ & $\boldsymbol{\lambda_1/\lambda_2}$ & \textbf{Gap $\boldsymbol{\delta}$} & \textbf{Erosion (\%)} & \textbf{Rotation (deg)} \\
\midrule
A6140 & 362 & 7.22 [3.98, 12.22] & 0.819 [0.441, 1.306] & --- & 6.5 [3.7, 12.1] \\
GA150 & 118 & 1.76 [1.12, 3.57] & 0.134 [0.026, 0.364] & 84 & 27.8 [9.1, 84.6] \\
GA250 & 121 & 1.79 [1.12, 3.45] & 0.118 [0.022, 0.301] & 86 & 32.8 [9.7, 85.5] \\
GA450 & 83 & 3.20 [1.44, 7.42] & 0.418 [0.119, 0.936] & 49 & 15.2 [5.8, 47.7] \\
\bottomrule
\end{tabular}
\caption{\textbf{Posterior summary of the channel-operator geometry in
\textit{Caenorhabditis elegans}.} Posterior median and 95\% credible
interval from 10,000 MCMCglmm draws. Erosion reports percent loss of
the spectral gap relative to the A6140 ancestor.}
\label{tab:celegans_summary}
\end{table}
